\documentclass[11pt,a4paper]{article}
\usepackage[utf8]{inputenc}
\usepackage[T1]{fontenc}
\usepackage[portuguese]{babel}
\usepackage{geometry}
\geometry{margin=2.5cm}
\usepackage{listings}
\usepackage{xcolor}
\usepackage{amsmath}
\usepackage{amssymb}
\usepackage{hyperref}
\usepackage{enumitem}
\usepackage{tcolorbox}
\usepackage{tabularx}
\usepackage{booktabs}

\definecolor{codebg}{rgb}{0.95,0.95,0.95}
\definecolor{codegreen}{rgb}{0.1,0.5,0.1}
\definecolor{codepurple}{rgb}{0.5,0.1,0.5}
\definecolor{codegray}{rgb}{0.5,0.5,0.5}

\lstdefinestyle{python}{
    backgroundcolor=\color{codebg},
    commentstyle=\color{codegray}\itshape,
    keywordstyle=\color{codepurple}\bfseries,
    stringstyle=\color{codegreen},
    basicstyle=\ttfamily\small,
    breaklines=true,
    frame=single,
    language=Python,
    showstringspaces=false,
    tabsize=4,
    literate={á}{{\'a}}1 {ã}{{\~a}}1 {é}{{\'e}}1 {í}{{\'i}}1 {ó}{{\'o}}1 {õ}{{\~o}}1 {ú}{{\'u}}1 {ç}{{\c{c}}}1 {Á}{{\'A}}1 {Ã}{{\~A}}1 {É}{{\'E}}1 {Í}{{\'I}}1 {Ó}{{\'O}}1 {Õ}{{\~O}}1 {Ú}{{\'U}}1 {Ç}{{\c{C}}}1
}
\lstset{style=python}

\newtcolorbox{objetivos}[1][]{
  colback=green!5!white,
  colframe=green!40!black,
  title=O que você vai aprender nesta página,
  fonttitle=\bfseries,
  #1
}

\title{\textbf{Data Analysis with Python 2024--2025}\\Parte 2: Python\\\large Conteúdo Teórico}
\author{Tradução e Adaptação: Universidade de Helsinque (MOOC.fi)}
\date{}

\begin{document}

\maketitle

\section*{Nota Importante}
Este material é uma tradução e adaptação baseada no curso \textit{Data Analysis with Python 2024-2025}, oferecido pela \textbf{Universidade de Helsinque (MOOC.fi)}.

\begin{objetivos}
\begin{itemize}
\item Compreender o funcionamento das expressões regulares em Python.
\item Aprender a processar arquivos em Python.
\item Entender os conceitos de objetos e classes.
\item Saber tratar exceções em um programa Python.
\end{itemize}
\end{objetivos}

\tableofcontents
\newpage

\section{Expressões regulares}

Já vimos que é possível perguntar a uma string \texttt{str} se ela começa
com um determinado prefixo usando \texttt{str.startswith('Apple')}. Se
quiséssemos verificar se ela começa com \texttt{"Apple"} \emph{ou}
\texttt{"apple"}, teríamos que chamar o método duas vezes. As expressões
regulares oferecem uma forma mais simples de resolver isso:
\texttt{re.match(r"[Aa]pple", str)}. A notação entre colchetes é um
exemplo da sintaxe especial das expressões regulares: ela indica que
qualquer um dos caracteres dentro dos colchetes é aceito -- neste caso,
\texttt{"A"} ou \texttt{"a"}. O restante dos caracteres em \texttt{"pple"}
funciona normalmente. A string \texttt{r"[Aa]pple"} é chamada de
\emph{padrão} (pattern).

Um exemplo um pouco mais elaborado verifica se a string começa com
\texttt{apple} ou \texttt{banana} (independentemente de maiúsculas ou
minúsculas na primeira letra):
\texttt{re.match(r"[Aa]pple|[Bb]anana", str)}. Aqui aparece um novo
caractere especial, a barra vertical \texttt{|}, que representa uma
alternativa. De cada lado da barra temos um subpadrão.

Um nome de variável válido em Python começa com uma letra ou um
sublinhado, e os caracteres seguintes também podem ser dígitos. Assim,
\texttt{\_oculto}, \texttt{L\_valor} e \texttt{A123\_} são nomes válidos,
mas \texttt{2abc} não é. O padrão de expressão regular para reconhecer
nomes de variáveis válidos seria:
\texttt{r"[A-Za-z\_][A-Za-z\_0-9]*\textbackslash Z"}. Aqui usamos uma
abreviação para intervalos de caracteres: \texttt{A-Z} significa todos os
caracteres de \texttt{A} até \texttt{Z}.

O primeiro caractere do nome é definido no primeiro par de colchetes; os
caracteres seguintes, no segundo par. O caractere especial \texttt{*}
significa ``qualquer quantidade (0, 1, 2, \dots) do subpadrão anterior''.
Por exemplo, o padrão \texttt{r"ba*"} aceita as strings \texttt{"b"},
\texttt{"ba"}, \texttt{"baa"}, \texttt{"baaa"} e assim por diante. A
notação especial \texttt{\textbackslash Z} indica o fim da string. Sem
ela, uma string como \texttt{abc-} também seria aceita, pois a função
\texttt{match} normalmente verifica apenas se a string \emph{começa} com
o padrão.

As notações especiais, como \texttt{\textbackslash Z}, podem gerar
confusão, já que os literais de string em Python também usam notações
parecidas: \texttt{\textbackslash n} representa uma quebra de linha,
\texttt{\textbackslash t} representa uma tabulação, e assim por diante.
Isso é resolvido usando as chamadas \emph{raw strings} (strings
``cruas''), indicadas por um \texttt{r} antes das aspas, por exemplo
\texttt{r"ab*\textbackslash Z"}. Em uma raw string, notações como
\texttt{\textbackslash n} e \texttt{\textbackslash t} não são
interpretadas. \textbf{É recomendável sempre usar raw strings ao definir
padrões de expressões regulares!}

\subsection{Padrões}

Um padrão representa um conjunto de strings -- potencialmente infinito --
que compartilham alguma característica em comum. Expressões regulares
(RE) são um tópico clássico da ciência da computação e são muito comuns
em tarefas de programação. Linguagens de script, como Python, lidam muito
bem com expressões regulares, permitindo processar textos de forma bem
sofisticada.

Em um padrão, caracteres normais (letras, números) representam a si
mesmos, a menos que sejam precedidos por uma barra invertida, o que pode
ativar um significado especial. Sinais de pontuação têm significado
especial, a não ser que sejam precedidos por \texttt{\textbackslash}, o
que remove esse significado especial. Use \texttt{\textbackslash
\textbackslash} para representar uma barra invertida literal.

\begin{center}
\begin{tabularx}{\textwidth}{lX}
\toprule
\textbf{Símbolo} & \textbf{Significado} \\
\midrule
\texttt{.} & Casa com qualquer caractere \\
\texttt{[...]} & Casa com qualquer caractere presente dentro dos colchetes \\
\texttt{[\^{}...]} & Casa com qualquer caractere que \emph{não} apareça após o acento circunflexo \\
\texttt{\^{}} & Casa com o início da string \\
\texttt{\$} & Casa com o final da string \\
\texttt{*} & Casa com zero ou mais ocorrências da RE anterior \\
\texttt{+} & Casa com uma ou mais ocorrências da RE anterior \\
\texttt{\{m,n\}} & Casa de m a n ocorrências da RE anterior \\
\texttt{?} & Casa com zero ou uma ocorrência da RE anterior \\
\bottomrule
\end{tabularx}
\end{center}

Já vimos que o caractere \texttt{|} denota alternativas. Por exemplo, o
padrão \texttt{r"Get (on|off|ready)"} casa com as strings
\texttt{"Get on"}, \texttt{"Get off"} e \texttt{"Get ready"}. Parênteses
podem ser usados para criar agrupamentos dentro de um padrão:
\texttt{r"(ab)+"} casa com \texttt{"ab"}, \texttt{"abab"},
\texttt{"ababab"}, etc. Esses grupos recebem também um número de
referência, a partir de 1, e podem ser referenciados dentro do próprio
padrão por meio de \emph{retro-referências} (\texttt{\textbackslash
número}). Por exemplo, podemos localizar padrões repetidos separados por
espaço com \texttt{r"([a-z]\{3,\}) \textbackslash 1 \textbackslash 1"}, o
que reconhece strings como \texttt{"aca aca aca"} ou
\texttt{"turn turn turn"}, mas não \texttt{"aca aba aca"} nem
\texttt{"ac ac ac"}.

Um acento circunflexo como primeiro caractere dentro de colchetes cria o
\emph{conjunto complementar} de caracteres:

\begin{center}
\begin{tabularx}{\textwidth}{lX}
\toprule
\textbf{Símbolo} & \textbf{Significado} \\
\midrule
\texttt{\textbackslash d} & equivalente a \texttt{[0-9]}, casa com um dígito \\
\texttt{\textbackslash D} & equivalente a \texttt{[\^{}0-9]}, casa com qualquer coisa exceto um dígito \\
\texttt{\textbackslash s} & casa com um caractere de espaço em branco (espaço, quebra de linha, tabulação, \dots) \\
\texttt{\textbackslash S} & casa com um caractere que não é espaço em branco \\
\texttt{\textbackslash w} & equivalente a \texttt{[a-zA-Z0-9\_]}, casa com um caractere alfanumérico \\
\texttt{\textbackslash W} & casa com um caractere não alfanumérico \\
\bottomrule
\end{tabularx}
\end{center}

Com essa notação, o exemplo de nome de variável pode ser reescrito de
forma mais curta como \texttt{r'[a-zA-Z\_]\textbackslash w*\textbackslash
Z'}.

Os padrões \texttt{\textbackslash A}, \texttt{\textbackslash b},
\texttt{\textbackslash B} e \texttt{\textbackslash Z} casam todos com uma
string vazia, mas em posições específicas. \texttt{\textbackslash A} e
\texttt{\textbackslash Z} reconhecem, respectivamente, o início e o fim
da string inteira -- diferentemente de \texttt{\^{}} e \texttt{\$}, que em
certos modos também podem casar após ou antes de uma quebra de linha. O
padrão \texttt{\textbackslash b} casa com o início ou o fim de uma
palavra; \texttt{\textbackslash B} faz o oposto.

\subsection{Funções de busca}

Até agora usamos apenas a função \texttt{re.match}, que tenta encontrar
uma correspondência no \emph{início} da string. A função \texttt{re.search}
permite localizar a correspondência em qualquer parte da string. Por
exemplo, \texttt{re.search(r'\textbackslash bback\textbackslash b', s)}
casa com \texttt{"back"}, \texttt{"a back, is a body part"} e
\texttt{"get back"}, mas não com \texttt{"backspace"} ou
\texttt{"comeback"}.

A função \texttt{re.search} encontra apenas a primeira ocorrência.
Podemos usar \texttt{re.findall} para encontrar todas as ocorrências.
Suponha que queremos encontrar todos os gerúndios (palavras terminadas em
\texttt{'ing'}) em uma string \texttt{s}:

\begin{lstlisting}
import re
s = "Doing things, going home, staying awake, sleeping later"
re.findall(r'\w+ing\b', s)
# ['Doing', 'going', 'staying', 'sleeping']
\end{lstlisting}

Para capturar todos os números inteiros de uma string, podemos usar
\texttt{re.findall(r'[+-]?\textbackslash d+', s)}:

\begin{lstlisting}
re.findall(r'[+-]?\d+', "23 + -24 = -1")
# ['23', '-24', '-1']
\end{lstlisting}

\subsubsection{Casamento guloso e não guloso}

Suponha que temos frases do tipo ``se \dots, então \dots'' e queremos
extrair apenas a condição. Uma primeira tentativa,
\texttt{re.findall(r'[Ii]f (.*), then', s)}, tende a capturar texto
demais, pois o padrão \texttt{.*} tenta casar o máximo possível de
caracteres -- esse comportamento é chamado de \emph{casamento guloso}
(\emph{greedy matching}).

Uma forma de resolver isso é impedir que o ponto final seja incluído no
trecho capturado, usando uma classe de caracteres complementar:
\texttt{[\^{}.]}, resultando em
\texttt{re.findall(r'[Ii]f ([\^{}.]*), then', s)}.

Outra solução é usar o \emph{casamento não guloso}. Os quantificadores
\texttt{+}, \texttt{*}, \texttt{?} e \texttt{\{m,n\}} possuem versões não
gulosas: \texttt{+?}, \texttt{*?}, \texttt{??} e \texttt{\{m,n\}?}. Essas
versões tentam usar o \emph{mínimo} de caracteres possível para que todo
o padrão ainda case com algum trecho da string:

\begin{lstlisting}
re.findall(r'[Ii]f (.*?), then', s)
\end{lstlisting}

\subsection{Funções do módulo \texttt{re}}

As funções mais comuns do módulo \texttt{re} são:

\begin{itemize}
\item \texttt{re.match(padrao, str)}
\item \texttt{re.search(padrao, str)}
\item \texttt{re.findall(padrao, str)}
\item \texttt{re.finditer(padrao, str)}
\item \texttt{re.sub(padrao, substituicao, str, count=0)}
\end{itemize}

As funções \texttt{match} e \texttt{search} retornam um \emph{objeto de
correspondência} (\emph{match object}), que descreve a ocorrência
encontrada. A função \texttt{findall} retorna uma lista com todas as
ocorrências (como strings). A função \texttt{finditer} funciona como
\texttt{findall}, mas retorna um iterador cujos itens são objetos de
correspondência. A função \texttt{sub} substitui todas as ocorrências do
padrão em \texttt{str} pela string de substituição e retorna a nova
string.

\begin{lstlisting}
import re
frase = "She goes where she wants to, she's a sheriff."
nova_frase = re.sub(r'\b[Ss]he\b', 'he', frase)
print(nova_frase)
# he goes where he wants to, he's a sheriff.
\end{lstlisting}

A função \texttt{sub} também aceita retro-referências para se referir ao
texto casado. As retro-referências \texttt{\textbackslash 1},
\texttt{\textbackslash 2} etc. referem-se, em ordem, aos grupos do
padrão:

\begin{lstlisting}
import re
texto = """He is a timelord.
He has a Tardis."""
novo_texto = re.sub(r'(\b[Hh]e\b)', r'\1 (The Doctor)', texto, 1)
print(novo_texto)
# He (The Doctor) is a timelord.
# He has a Tardis.
\end{lstlisting}

\subsection{O objeto de correspondência (\emph{match object})}

As funções \texttt{match}, \texttt{search} e \texttt{finditer} usam
objetos de correspondência para descrever a ocorrência encontrada. O
método \texttt{groups()} de um objeto de correspondência retorna uma
tupla com todas as substrings casadas pelos grupos do padrão. Cada par de
parênteses no padrão cria um novo grupo, referenciado pelos índices 1, 2,
\dots{} O grupo 0 é especial: representa o casamento produzido pelo
padrão inteiro.

\begin{lstlisting}
mo = re.search(r'\d+ (\d+) \d+ (\d+)', 'first 123 45 67 890 last')
mo.groups()   # ('45', '890')
mo.group(1)   # '45'
mo.group(0)   # '123 45 67 890'
\end{lstlisting}

Além de acessar as strings casadas pelo padrão e seus grupos, também é
possível obter os índices correspondentes na string original:

\begin{itemize}
\item \texttt{start(gid=0)} e \texttt{end(gid=0)} retornam, respectivamente, os índices de início e fim do grupo \texttt{gid};
\item \texttt{span(gid)} retorna o par (início, fim).
\end{itemize}

O objeto de correspondência também pode ser usado como um valor booleano:

\begin{lstlisting}
mo = re.search(...)
if mo:
    # faca algo
\end{lstlisting}

Ou, de forma explícita, \texttt{encontrado = bool(mo)}.

\subsection{Outras informações úteis}

Quando o mesmo padrão é usado em várias chamadas, pode ser interessante
\emph{pré-compilá-lo} por questões de eficiência, usando a função
\texttt{compile(padrao, flags=0)} do módulo \texttt{re}. Essa função
retorna um chamado \emph{objeto RE}, que possui versões dos métodos do
módulo \texttt{re}, com a diferença de que o padrão já está armazenado no
objeto (não é mais o primeiro parâmetro).

Detalhes do casamento podem ser ajustados por \emph{flags} opcionais,
fornecidas dentro do próprio padrão ou como segundo parâmetro da função
\texttt{compile}:

\begin{center}
\begin{tabular}{ll}
\toprule
\textbf{Notação no padrão} & \textbf{Flag} \\
\midrule
\texttt{(?i)} & \texttt{re.IGNORECASE} \\
\texttt{(?m)} & \texttt{re.MULTILINE} \\
\texttt{(?s)} & \texttt{re.DOTALL} \\
\bottomrule
\end{tabular}
\end{center}

A flag \texttt{IGNORECASE} faz com que letras maiúsculas e minúsculas
sejam tratadas como equivalentes. A flag \texttt{MULTILINE} faz com que
\texttt{\^{}} e \texttt{\$} casem com o início e o fim de \emph{cada
linha}, além do início e fim da string inteira -- isso é o que diferencia
\texttt{\textbackslash A} de \texttt{\^{}}, e \texttt{\textbackslash Z} de
\texttt{\$}. A flag \texttt{DOTALL} faz com que o ponto (\texttt{.})
também aceite o caractere de quebra de linha. Múltiplas flags podem ser
combinadas com \texttt{|}: \texttt{re.compile(padrao, re.MULTILINE |
re.DOTALL)}.

\section{Processamento básico de arquivos}

Um arquivo pode ser aberto com a função \texttt{open}. A chamada
\texttt{open(nome\_do\_arquivo, mode="r")} retorna um objeto de arquivo,
usado para referenciar um arquivo em disco. Para ler ou escrever, usamos
os métodos \texttt{read} e \texttt{write} desse objeto. Ao terminar de
usá-lo, deve-se chamar o método \texttt{close}.

O parâmetro \texttt{mode} controla que tipo de operação pode ser feita no
arquivo:

\begin{center}
\begin{tabularx}{\textwidth}{lX}
\toprule
\textbf{Modo} & \textbf{Descrição} \\
\midrule
\texttt{r} & somente leitura; o arquivo precisa existir \\
\texttt{w} & somente escrita; cria o arquivo, ou sobrescreve se já existir \\
\texttt{a} & somente escrita; a escrita sempre ocorre no final do arquivo (append) \\
\texttt{r+} & leitura e escrita; o arquivo precisa existir \\
\texttt{w+} & leitura e escrita; cria, ou sobrescreve se já existir \\
\texttt{a+} & leitura e escrita; a escrita ocorre no final do arquivo \\
\bottomrule
\end{tabularx}
\end{center}

Ao final da string de modo, pode-se acrescentar a letra \texttt{t} (modo
texto) ou \texttt{b} (modo binário); se omitida, o padrão é o modo texto.

No modo binário, o conteúdo do arquivo não é interpretado de forma
alguma: os métodos \texttt{read} e \texttt{write} lidam com \emph{bytes}
(um byte tem 8 bits e pode representar um número entre 0 e 255).

No modo texto, duas conversões ocorrem automaticamente:

\begin{itemize}
\item No Windows, o fim de linha é codificado por dois caracteres; ao ler
  o arquivo, esses dois caracteres são convertidos para \texttt{'\textbackslash
  n'}, e o processo inverso ocorre ao escrever.
\item Cada caractere é codificado no arquivo como um ou mais bytes. Uma
  codificação muito usada é a UTF-8. Nessa codificação, por exemplo, o
  caractere \texttt{'ã'} é representado por dois bytes.
\end{itemize}

Neste curso trabalharemos apenas com arquivos de texto, então o modo
texto padrão é suficiente. Ainda assim, às vezes é necessário especificar
explicitamente a codificação de um arquivo, caso não seja UTF-8.

\subsection{Métodos comuns do objeto de arquivo}

\begin{itemize}
\item \texttt{read(tamanho)}: lê \texttt{tamanho} caracteres/bytes como string.
\item \texttt{write(string)}: escreve uma string/bytes no arquivo.
\item \texttt{readline()}: lê uma linha, incluindo o caractere de quebra de linha final.
\item \texttt{readlines()}: retorna uma lista com todas as linhas do arquivo.
\item \texttt{writelines()}: escreve uma lista de linhas no arquivo.
\item \texttt{flush()}: tenta garantir que as alterações sejam gravadas em disco imediatamente.
\end{itemize}

\begin{lstlisting}
f = open("caderno.ipynb", "r")   # abre um arquivo de texto
for i in range(5):               # le as cinco primeiras linhas
    linha = f.readline()
    print(f"Linha {i}: {linha}", end="")
f.close()
\end{lstlisting}

É fácil esquecer de fechar um arquivo. Uma solução é usar um
\emph{gerenciador de contexto}, criado com a instrução \texttt{with}: ao
sair do bloco indentado, o arquivo é fechado automaticamente.

\begin{lstlisting}
with open("caderno.ipynb", "r") as f:   # o arquivo sera fechado
    for i in range(5):                  # automaticamente ao sair do bloco
        linha = f.readline()
        print(f"Linha {i}: {linha}", end="")
\end{lstlisting}

O objeto de arquivo é iterável, ou seja, podemos percorrer as linhas do
arquivo usando um laço \texttt{for}:

\begin{lstlisting}
maior_tamanho = 0
with open("caderno.ipynb", "r") as f:
    for linha in f:                     # percorre todas as linhas do arquivo
        if len(linha) > maior_tamanho:
            maior_tamanho = len(linha)
print(f"A maior linha deste arquivo tem tamanho {maior_tamanho}")
\end{lstlisting}

\subsection{Objetos padrão de entrada e saída}

O Python mantém três objetos de arquivo automaticamente abertos:

\begin{itemize}
\item \texttt{sys.stdin} -- entrada padrão;
\item \texttt{sys.stdout} -- saída padrão;
\item \texttt{sys.stderr} -- saída padrão de erro.
\end{itemize}

Para ler uma linha digitada pelo usuário, usa-se
\texttt{sys.stdin.readline()}. Para escrever uma linha na tela, usa-se
\texttt{sys.stdout.write(linha)}. A saída de erro deve ser usada apenas
para mensagens de erro, ainda que, muitas vezes, seu destino final seja o
mesmo da saída padrão.

A função \texttt{print} usa o arquivo \texttt{sys.stdout}, e a função
\texttt{input} usa o arquivo \texttt{sys.stdin}. Esses objetos podem ter
seu destino redirecionado -- é comum, por exemplo, redirecionar
\texttt{stderr} para um arquivo de log.

\subsection{O módulo \texttt{sys}}

Além dos três objetos de arquivo padrão, o módulo \texttt{sys} tem outros
atributos úteis. \texttt{sys.path} é a lista de pastas em que o Python
procura por módulos a serem importados. \texttt{sys.argv} contém os
chamados \emph{parâmetros de linha de comando}: ao executar
\texttt{python3 programa.py param1 param2 \dots} no terminal, o nome do
programa fica em \texttt{sys.argv[0]}, e os demais parâmetros ficam nas
posições seguintes da lista.

A função \texttt{sys.exit} pode ser usada para encerrar imediatamente o
programa. O parâmetro inteiro fornecido é o valor de retorno do programa
-- por convenção, \texttt{0} indica sucesso, e qualquer valor diferente
de zero indica que ocorreu um erro.

\section{Objetos e classes}

Python é uma linguagem orientada a objetos, como Java e C++, mas,
diferentemente de Java, não obriga o uso de classes, herança e métodos:
também é possível programar de forma estrutural, com funções e módulos.

Todo valor em Python é um objeto. Objetos combinam dados e as funções que
manipulam esses dados -- essa combinação é chamada de
\emph{encapsulamento}. Os itens de dados e as funções de um objeto são
chamados de \emph{atributos}; em particular, os atributos que são funções
são chamados de \emph{métodos}. Por exemplo, o operador \texttt{+} em
números inteiros invoca um método dos inteiros, e o operador \texttt{+}
em strings invoca um método das strings.

Funções, módulos, métodos, classes etc. são todos \emph{objetos de
primeira classe} em Python, ou seja, podem ser:

\begin{itemize}
\item armazenados em um contêiner;
\item passados como parâmetro para uma função;
\item retornados por uma função;
\item associados a uma variável.
\end{itemize}

Acessamos um atributo de um objeto usando o \emph{operador ponto}:
\texttt{objeto.atributo}. Por exemplo, se \texttt{L} é uma lista,
referenciamos o método \texttt{append} com \texttt{L.append}, e podemos
chamá-lo assim: \texttt{L.append(4)}. Como módulos também são objetos em
Python, a expressão \texttt{math.pi} nada mais é do que o acesso ao
atributo de dado \texttt{pi} do objeto módulo \texttt{math}.

Números como \texttt{2} e \texttt{100} são instâncias do tipo
\texttt{int}. Da mesma forma, \texttt{"ola"} é uma instância do tipo
\texttt{str}. Quando escrevemos \texttt{s = set()}, estamos, na verdade,
criando uma nova instância do tipo \texttt{set} e associando esse objeto
à variável \texttt{s}.

O usuário pode definir seus próprios tipos de dado -- as chamadas
\emph{classes}. Ao chamar uma classe como se fosse uma função, obtemos um
novo objeto \emph{instância} desse tipo. Classes podem ser vistas como
``receitas'' para criar objetos.

\begin{lstlisting}
class MinhaClasse(object):
    """String de documentacao da classe"""

    def __init__(self, param1, param2):
        "Inicializa uma instancia de MinhaClasse"
        self.b = param1     # cria um atributo de instancia
        c = param2           # cria uma variavel local da funcao
        # demais instrucoes ...

    def f(self, param1):
        """Este e um metodo da classe"""
        # algumas instrucoes

    a = 1   # isto cria um atributo de classe
\end{lstlisting}

A definição de classe começa com a palavra-chave \texttt{class}. Nela
damos um nome ao novo tipo e, entre parênteses, listamos as classes-base.
O bloco indentado seguinte é o \emph{corpo da classe}. Depois de todo o
corpo ser lido, um novo tipo é criado -- note que nenhuma instância é
criada ainda. Todos os atributos e métodos da classe são definidos no
corpo da classe.

A classe do exemplo tem dois métodos: \texttt{\_\_init\_\_} e \texttt{f}.
O primeiro parâmetro deles é especial: \texttt{self} refere-se à
instância da classe sobre a qual o método está sendo chamado. Por
exemplo, se criarmos \texttt{a = MinhaClasse(...)}, então, na chamada
\texttt{a.f(param1)}, \texttt{self} representa \texttt{a}. O conceito de
\texttt{self} é equivalente ao de \texttt{this} em C++ ou Java.

O método \texttt{\_\_init\_\_} realiza a inicialização quando uma
instância é criada. Ao instanciar com \texttt{i = MinhaClasse(2,3)}, os
parâmetros \texttt{param1} e \texttt{param2} recebem os valores 2 e 3,
respectivamente. Com a instância \texttt{i} criada, podemos chamar seu
método \texttt{f} com o operador ponto: \texttt{i.f(1)}.

Há diferenças em como uma atribuição dentro do corpo da classe cria
variáveis. O atributo \texttt{a} é definido no nível da classe e é comum
a todas as instâncias de \texttt{MinhaClasse}. A variável \texttt{c} é
local à função \texttt{\_\_init\_\_} e não pode ser usada fora dela. Já o
atributo \texttt{b} é específico de cada instância. Assim, para objetos
\texttt{x = MinhaClasse(1,0)} e \texttt{y = MinhaClasse(2,0)}, temos
\texttt{x.b != y.b}, mas \texttt{x.a == y.a}.

Todo método de uma classe recebe obrigatoriamente um primeiro parâmetro
que se refere à instância na qual o método foi chamado, chamado
convencionalmente de \texttt{self}. Para acessar um atributo de classe a
partir de um método, use a forma completa \texttt{MinhaClasse.a}. Métodos
cujo nome começa e termina com dois sublinhados são chamados de
\emph{métodos especiais}; \texttt{\_\_init\_\_} é um deles.

\subsection{Instâncias}

Criamos instâncias chamando uma classe como se fosse uma função:
\texttt{i = NomeDaClasse(...)}. Os parâmetros fornecidos na chamada são
passados para a função \texttt{\_\_init\_\_}, onde os atributos
específicos da instância podem ser criados. Se \texttt{\_\_init\_\_}
estiver ausente, é possível criar uma instância sem fornecer nenhum
parâmetro -- nesse caso, ela não terá atributos. Mais tarde, é possível
(re)associar atributos com a atribuição \texttt{instancia.atributo =
novo\_valor}.

Se o atributo não existia antes, ele é adicionado à instância com o valor
atribuído. Em Python é realmente possível adicionar ou remover atributos
de uma instância já existente, porque os nomes dos atributos e seus
respectivos valores são armazenados em um dicionário, que é, ele mesmo,
um atributo da instância chamado \texttt{\_\_dict\_\_}. Outro atributo
padrão, além de \texttt{\_\_dict\_\_}, é \texttt{\_\_class\_\_}, que
armazena a classe da instância -- ou seja, o tipo do objeto.

\subsection{Busca de atributos}

Suponha que \texttt{x} é uma instância da classe \texttt{X}, e queremos
ler um atributo \texttt{x.a}. A busca ocorre em três etapas:

\begin{enumerate}
\item verifica-se primeiro se \texttt{a} é um atributo da instância \texttt{x};
\item caso não seja, verifica-se se \texttt{a} é um atributo de classe de \texttt{X};
\item caso ainda não seja, as classes-base de \texttt{X} são verificadas.
\end{enumerate}

Já para \emph{associar} um atributo, as regras são mais simples:
\texttt{x.a = valor} define o atributo de instância, e \texttt{X.a =
valor} define o atributo de classe. Se uma classe-base, a própria classe
\texttt{X} e a instância \texttt{x} tiverem, cada uma, um atributo chamado
\texttt{a}, então \texttt{x.a} esconde \texttt{X.a}, que por sua vez
esconde o atributo da classe-base.

\subsection{Herança}

A herança nos permite reaproveitar o código de uma classe já existente
\texttt{B} na criação de uma nova classe \texttt{C}. Relembrando: ao
procurar um atributo, a busca começa no dicionário da instância e
continua pelos atributos de classe; se ambos falharem, o atributo é
buscado nas classes-base e, recursivamente, nas classes-base delas.

Assim, pode parecer que estamos acessando um atributo da classe
\texttt{C}, quando na verdade estamos acessando um atributo herdado de
sua classe-base \texttt{B}. Nesse caso, dizemos que \texttt{C}
\emph{herda} o atributo de sua classe-base \texttt{B}. Se houver
atributos com o mesmo nome tanto na classe quanto na sua classe-base, o
atributo da classe-base fica oculto -- dizemos que a classe \texttt{C}
\emph{sobrescreve} (\emph{override}) o atributo da classe-base
\texttt{B}. Terminologia: \texttt{B} é a classe-base (superclasse) e
\texttt{C} é a classe derivada (subclasse).

\begin{lstlisting}
class B(object):
    def f(self):
        print("Executando B.f")
    def g(self):
        print("Executando B.g")

class C(B):
    def g(self):
        print("Executando C.g")

x = C()
x.f()   # herdado de B
x.g()   # sobrescrito por C
# Saida:
# Executando B.f
# Executando C.g
\end{lstlisting}

A relação de herança entre duas classes \texttt{B} e \texttt{C} pode ser
testada com a função \texttt{issubclass}:
\texttt{issubclass(C, B) == True}, mas \texttt{issubclass(B, C) ==
False}. A função \texttt{isinstance(obj, cls)} permite verificar se uma
instância tem tipo \texttt{cls} ou algum ancestral do tipo \texttt{cls}.
Criando \texttt{x = C()} e \texttt{y = B()}, temos
\texttt{isinstance(x, B) == isinstance(x, C) == isinstance(y, B) ==
True}, mas \texttt{isinstance(y, C) == False}.

A classe \texttt{object} deve ser uma classe-base ou ancestral de toda
classe em Python: para toda instância \texttt{x},
\texttt{isinstance(x, object) == True}.

Ao derivar de uma classe existente, podemos modificar e/ou estender seu
comportamento sem alterar a classe original. Por exemplo, se quisermos
adicionar um método à classe \texttt{list}, podemos usar herança e
codificar apenas a parte que muda. Outro uso da herança é criar
hierarquias conceituais -- como veremos na hierarquia de exceções do
Python. Um terceiro uso é criar interfaces: várias classes podem
compartilhar a mesma interface (isto é, oferecer os mesmos atributos),
mas ter comportamentos ou implementações muito diferentes, o que permite
trocar parte do programa com mínimas alterações no restante do código.

Se, na definição de um método \texttt{C.f}, precisarmos chamar o método
correspondente da classe-base \texttt{B}, podemos usar a chamada completa
\texttt{B.f(...)}. Isso é chamado de \emph{delegação}, útil, por exemplo,
para chamar o \texttt{\_\_init\_\_} da classe-base a partir do
\texttt{\_\_init\_\_} da classe derivada, inicializando os atributos da
classe-base.

\subsection{Métodos especiais}

Já conhecemos um método especial, o \texttt{\_\_init\_\_}, que define os
valores iniciais dos atributos de instância. A seguir, apresentamos os
principais métodos especiais de propósito geral, executados quando certas
operações são realizadas sobre os objetos.

Sendo \texttt{C} uma classe e \texttt{x}, \texttt{y} suas instâncias:
\texttt{\_\_hash\_\_} retorna um valor inteiro, com o requisito de que
\texttt{x == y} implica \texttt{x.\_\_hash\_\_() == y.\_\_hash\_\_()}. O
valor é usado ao armazenar objetos em dicionários e conjuntos; as
instâncias devem ser imutáveis. Uma classe com o método
\texttt{\_\_call\_\_} torna suas instâncias ``chamáveis'': a chamada
\texttt{x(a, b, \dots)} invocará esse método especial com os parâmetros
fornecidos. O método \texttt{\_\_del\_\_} é chamado quando a instância
correspondente é destruída. O método \texttt{\_\_new\_\_} controla a
criação de novas instâncias, podendo ser usado, por exemplo, para criar
classes que possuem apenas uma instância.

O método \texttt{\_\_str\_\_} é chamado quando \texttt{print} precisa
exibir o valor de uma instância -- deve retornar uma string. O método
\texttt{\_\_repr\_\_} é chamado quando o interpretador interativo exibe o
valor de uma expressão avaliada; retorna uma representação ``canônica'',
que (ao menos em teoria) poderia ser usada para recriar o objeto
original. Os métodos especiais \texttt{\_\_eq\_\_}, \texttt{\_\_ge\_\_},
\texttt{\_\_gt\_\_}, \texttt{\_\_le\_\_}, \texttt{\_\_lt\_\_} e
\texttt{\_\_ne\_\_} são chamados quando os operadores correspondentes
\texttt{==}, \texttt{>=}, \texttt{>}, \texttt{<=}, \texttt{<} e
\texttt{!=} são usados.

Para que instâncias de uma classe suportem operações numéricas (\texttt{+},
\texttt{-}, \texttt{*}, \texttt{/} etc.), é preciso definir os métodos
especiais correspondentes. Por exemplo, a expressão \texttt{x + y}
resulta na chamada \texttt{x.\_\_add\_\_(y)}, que deve retornar o
resultado da operação.

\begin{center}
\begin{tabular}{ll}
\toprule
\textbf{Método} & \textbf{Operação} \\
\midrule
\texttt{\_\_add\_\_} & adição (+) \\
\texttt{\_\_sub\_\_} & subtração (-) \\
\texttt{\_\_mul\_\_} & multiplicação (*) \\
\texttt{\_\_truediv\_\_} & divisão (/) \\
\texttt{\_\_floordiv\_\_} & divisão inteira (//) \\
\bottomrule
\end{tabular}
\end{center}

As atribuições compostas correspondentes (\texttt{+=}, \texttt{-=},
\texttt{*=}, \texttt{/=}) usam os métodos especiais \texttt{\_\_iadd\_\_},
\texttt{\_\_isub\_\_}, \texttt{\_\_imul\_\_} e \texttt{\_\_itruediv\_\_}.
As funções de conversão \texttt{complex()}, \texttt{float()} e
\texttt{int()} chamam, respectivamente, \texttt{\_\_complex\_\_},
\texttt{\_\_float\_\_} e \texttt{\_\_int\_\_}.

Além dos métodos habituais dos contêineres, como \texttt{append} de uma
lista, várias operações são resolvidas por chamadas a métodos especiais
da classe do contêiner. O teste de pertencimento \texttt{x in c}
corresponde à chamada \texttt{c.\_\_contains\_\_(x)}. A exclusão de um
elemento com \texttt{del c[chave]} corresponde a
\texttt{c.\_\_delitem\_\_(chave)}. A leitura de um item com
\texttt{c[chave]} corresponde a \texttt{c.\_\_getitem\_\_(chave)}, e a
atribuição \texttt{c[chave] = valor} corresponde a
\texttt{c.\_\_setitem\_\_(chave, valor)}. O número de elementos de um
contêiner \texttt{c} é obtido com \texttt{len(c)}, que na verdade chama
\texttt{c.\_\_len\_\_()}. A chamada \texttt{iter(c)} invoca o método
especial \texttt{\_\_iter\_\_}.

\section{Exceções}

Quando ocorre um erro, o que podemos fazer?

\begin{itemize}
\item Exibir uma mensagem de erro;
\item Interromper a execução do programa;
\item Sinalizar o erro retornando um valor especial, como \texttt{-1} ou \texttt{None};
\item Ignorar o erro;
\item \dots
\end{itemize}

Essas soluções tendem a misturar a \emph{indicação} do problema com a
\emph{reação} a esse problema. O comportamento do programa diante de uma
situação de erro fica fixo na implementação da função, e a instância que
chamou a função nem sempre sabe o que fazer diante da situação.

A maioria das linguagens modernas possui um sistema chamado
\emph{tratamento de exceções}, que separa o reconhecimento de erros do
seu tratamento propriamente dito. Sinalizamos um erro ou situação
anômala \emph{lançando} (\emph{raising}) uma exceção, o que em Python é
feito com a instrução \texttt{raise}:

\begin{itemize}
\item \texttt{raise instancia}
\item \texttt{raise classe\_de\_excecao[, expressao]}
\end{itemize}

Na segunda forma, se a expressão existir, ela é uma tupla de parâmetros
passados para a classe de exceção.

As funções da biblioteca padrão do Python lançam exceções em situações de
erro. Às vezes, uma exceção nem representa exatamente um erro -- por
exemplo, quando um iterador se esgota, ele sinaliza isso lançando a
exceção \texttt{StopIteration}. Outro exemplo de exceção pouco ``grave''
é \texttt{Warning}.

A forma geral da instrução de captura de exceções é a seguinte:

\begin{lstlisting}
try:
    # instrucoes que podem causar excecoes
except (NomeExcecao1, NomeExcecao2, ...):
    # tratamento das excecoes
else:
    # executado se o bloco try nao causou nenhuma excecao
finally:
    # sempre executado -- codigo de limpeza
\end{lstlisting}

Em geral, apenas as partes \texttt{try} e \texttt{except} são necessárias.

\begin{lstlisting}
L = [1, 2, 3]
try:
    print(L[3])
except IndexError:
    print("Este indice nao existe")
# Saida: Este indice nao existe
\end{lstlisting}

\begin{lstlisting}
def calcula_media(L):
    n = len(L)
    s = sum(L)
    return float(s) / n   # o erro e detectado aqui!

minha_lista = []
while True:
    try:
        x = float(input("Digite um numero (nao numerico encerra): "))
        minha_lista.append(x)
    except ValueError:
        break

try:
    media = calcula_media(minha_lista)
    print("A media e", media)
except ZeroDivisionError:
    # e o erro e tratado aqui
    if len(minha_lista) == 0:
        print("Tentou calcular a media de uma lista vazia de numeros")
    else:
        print("Algo estranho aconteceu")
\end{lstlisting}

\subsection{A hierarquia de exceções}

Em Python, exceções são objetos, como qualquer outro valor, instanciados
a partir de classes de exceção. As classes de exceção formam hierarquias
naturais:

\begin{itemize}
\item novas classes de exceção podem ser criadas herdando de classes de exceção já existentes e estendendo-as;
\item a raiz dessa hierarquia é a classe \texttt{Exception};
\item o Python define diversas classes-base para servir de ponto de partida, além de várias classes de exceção prontas para uso.
\end{itemize}

Algumas das classes mais comuns e seus relacionamentos:

\begin{center}
\small
\begin{tabular}{ll}
\toprule
\textbf{Classe-base} & \textbf{Algumas classes derivadas} \\
\midrule
\texttt{Exception} & \texttt{SystemExit}, \texttt{StandardError}, \texttt{StopIteration}, \texttt{Warning} \\
\texttt{StandardError} & \texttt{NameError}, \texttt{SyntaxError}, \texttt{ArithmeticError}, \texttt{LookupError}, \texttt{EnvironmentError}, \dots \\
\texttt{ArithmeticError} & \texttt{FloatingPointError}, \texttt{OverflowError}, \texttt{ZeroDivisionError} \\
\texttt{LookupError} & \texttt{IndexError}, \texttt{KeyError} \\
\texttt{EnvironmentError} & \texttt{IOError}, \texttt{OSError} \\
\bottomrule
\end{tabular}
\end{center}

\subsection{Especificações de exceção genéricas demais}

A hierarquia de exceções permite capturar múltiplas exceções semelhantes
capturando sua classe-base comum. Esse recurso deve ser usado com
cuidado: uma especificação de exceção genérica demais, como
\texttt{except Exception:}, pode esconder a verdadeira causa de um erro.

\begin{lstlisting}
import sys
s = input("Digite um numero: ")
s = s[:-1]   # remove o caractere \n do final
try:
    x = int(s)
    sys.stdout.write(f"Voce digitou {x}\n")
except Exception:
    print("Voce nao digitou um numero")
\end{lstlisting}

Se o usuário não digitar uma string que representa um inteiro, a função
\texttt{int} lança \texttt{ValueError}. Em vez de capturar
\texttt{ValueError} especificamente, capturamos a raiz da hierarquia,
\texttt{Exception}, o que acaba capturando \emph{qualquer} exceção
possível -- inclusive um eventual erro de digitação no próprio código,
que passaria despercebido. Trocar a especificação de \texttt{Exception}
para \texttt{ValueError} deixaria esse tipo de problema visível.

\subsection{A política de tratamento de erros do Python}

O Python usa uma abordagem diferente da de muitas outras linguagens no
que diz respeito à checagem de erros. Em vez de verificar previamente se
todas as entradas têm o tipo correto e se o conteúdo das variáveis faz
sentido para determinada operação, o Python primeiro tenta executar a
operação e depois verifica se ela causou alguma exceção. Isso está
relacionado ao conceito de \emph{duck typing}: uma função funciona para
um conjunto de entradas se todas as operações no corpo da função fizerem
sentido para essas entradas. É por isso que os parâmetros das funções
normalmente não têm um tipo especificado.

\section{Sequências, iteráveis e geradores: revisão}

Em termos simples, um contêiner é \emph{iterável} se pudermos percorrer
todos os seus elementos usando um laço \texttt{for}. Todas as sequências
são iteráveis, mas existem outros objetos iteráveis além delas -- e
podemos até criar nossos próprios tipos iteráveis. Para isso, a classe
precisa ter um método especial \texttt{\_\_iter\_\_}, que retorna um
\emph{iterador} para o contêiner. Um iterador é um objeto que possui o
método \texttt{\_\_next\_\_}, responsável por retornar o próximo elemento
do contêiner.

\begin{lstlisting}
class IteradorDiaDaSemana(object):
    """Iterador sobre os dias da semana."""

    def __init__(self):
        self.i = 0   # comeca na segunda-feira
        self.dias = ("Segunda", "Terca", "Quarta", "Quinta",
                     "Sexta", "Sabado", "Domingo")

    def __iter__(self):
        # se este objeto fosse um container, este metodo
        # retornaria o iterador sobre os elementos do container.
        # como este objeto ja e um iterador, retornamos self.
        return self

    def __next__(self):
        if self.i == 7:
            raise StopIteration   # sinaliza que todos os dias ja foram percorridos
        else:
            dia = self.dias[self.i]
            self.i += 1
            return dia

for dia in IteradorDiaDaSemana():
    print(dia)
\end{lstlisting}

Podemos verificar se \texttt{IteradorDiaDaSemana} é do tipo
\texttt{Sequence}:

\begin{lstlisting}
from collections import abc   # obtem as classes-base abstratas

containers = ["abc", [1, 2, 3], (4, 5), IteradorDiaDaSemana()]
for c in containers:
    if isinstance(c, abc.Sequence):
        print(c, "e uma sequencia")
    else:
        print(c, "nao e uma sequencia")
\end{lstlisting}

O iterador não é uma sequência, pois, por exemplo, não é possível indexá-lo
com colchetes (\texttt{[]}), mas ele é, sim, um iterável:
\texttt{isinstance(IteradorDiaDaSemana(), abc.Iterable)} retorna
\texttt{True}.

É possível criar iteradores manualmente, mas a sintaxe fica um tanto
complicada. Uma alternativa mais simples são os \emph{geradores}. Um
gerador é uma função que contém uma instrução \texttt{yield}. Note a
diferença entre geradores e expressões geradoras (vistas anteriormente no
curso): ambos produzem iteráveis, mas de formas diferentes.

\begin{lstlisting}
def minhas_datas(dia=1, mes=1):
    """Gera datas a partir da data informada."""
    duracoes = (31, 28, 31, 30, 31, 30, 31, 31, 30, 31, 30, 31)
    primeiro_dia = dia
    for m in range(mes, 13):
        for d in range(primeiro_dia, duracoes[m - 1] + 1):
            yield (d, m)
        primeiro_dia = 1

# cria o gerador chamando a funcao:
gen = minhas_datas(26, 2)   # comeca em 26 de fevereiro

for i, (dia, mes) in enumerate(gen):
    if i == 5:
        break   # imprime apenas as cinco primeiras datas do gerador
    print(f"Indice {i}, dia {dia}, mes {mes}")
\end{lstlisting}

Não seria possível escrever esse mesmo iterável usando uma expressão
geradora, e seria bastante trabalhoso escrevê-lo explicitamente como um
iterador, como fizemos com \texttt{IteradorDiaDaSemana}.

\section*{Referências e créditos}
\addcontentsline{toc}{section}{Referências e créditos}

Este material é uma tradução e adaptação, para a língua portuguesa, do
conteúdo do curso \emph{Data Analysis with Python 2024--2025}, Capítulo
2 -- ``Python (Part 2)'', originalmente disponibilizado em inglês pela
plataforma MOOC.fi:

\begin{quote}
\url{https://courses.mooc.fi/org/uh-cs/courses/data-analysis-with-python-2024-2025/chapter-2/python-continues}
\end{quote}

\end{document}
